\documentclass[11pt]{article}
\usepackage[margin=1in]{geometry}
\usepackage{amsmath,amssymb,bm}
\usepackage{hyperref}
\newcommand{\E}{\mathbb{E}}
\newcommand{\MI}{\mathrm{MI}}
\newcommand{\soft}{\mathrm{soft}}
\newcommand{\Potts}{\textsc{Potts}}
\newcommand{\CTMC}{\textsc{ctmc}}

\title{A null-spiked, $L_1$-sparse prior for coevolutionary \Potts\ mixtures}
\author{Companion derivation for \texttt{paper2b-symmetry.tex}}
\date{}

\begin{document}
\maketitle

\noindent\emph{Status: derivation, ahead of implementation. Notation follows
\texttt{paper2b-symmetry.tex}: a mixture of $K$ coevolutionary ``interaction classes''
fit to a corpus of contacting column pairs, each class a reversible pair \CTMC, with a
shared single-site exchangeability $S$ and a $\Gamma$+I rate grid orthogonal to the class
label. The empirical finding this prior is built for: a mixture of coupled
(\Potts) classes beats a mixture of null products-of-marginals by only $\sim$$0.005$--$0.008$
nats/count, so coupling is a small, class-dependent effect. We want a prior that lets each
class decide whether it needs a coupling term at all, and, if so, keeps only the pairs the
data support --- turning ``how much coupling survives as phylogenetic signal'' into an
inferred model-selection quantity rather than a blanket assertion.}

\section{The class model and its null}
\label{sec:model}

Write a class-$k$ pair stationary in \Potts\ form on the $A{=}20$ amino acids,
\begin{equation}
\label{eq:stationary}
\pi_k(a,b)=\frac{1}{Z_k}\exp\!\big(h^{(k)}_a+h^{(k)}_b+J^{(k)}_{ab}\big),\qquad
J^{(k)}=J^{(k)\top},\quad J^{(k)}_{aa}=0,
\end{equation}
with \emph{fields} $h^{(k)}\in\mathbb R^{A}$ (single-site composition) and a symmetric,
hollow \emph{coupling} $J^{(k)}$. The class reversible generator is built from the shared
exchangeability $S$ by the Metropolis-$\sqrt{\cdot}$ rule, e.g.\ for a site-1 move
$(a,c)\!\to\!(b,c)$,
\begin{equation}
\label{eq:generator}
Q_k\big[(a,c)\!\to\!(b,c)\big]=S_{ab}\,\sqrt{\pi_k(b,c)/\pi_k(a,c)}
   =S_{ab}\,\exp\!\tfrac12\!\big[(h_b-h_a)+(J_{bc}-J_{ac})\big],
\end{equation}
and symmetrically for site 2; simultaneous moves have rate $0$ (single-transition class).
The partition function $Z_k$ \emph{cancels} from every transition rate, so the
transition likelihood depends on $J^{(k)}$ only through \emph{differences}
$J_{bc}-J_{ac}$ --- the same cancellation that lets the coevolution ELBO reduce to pairwise
bridge statistics (paper2b, \S\,fixed-bridge). The equilibrium enters the likelihood only
through the dwell/exposure term (below), never through $\log Z_k$.

\paragraph{The null is $J^{(k)}=0$.} Setting the coupling to zero collapses
\eqref{eq:stationary} to a rank-one additive energy,
\begin{equation}
\pi_k(a,b)\big|_{J=0}=\theta^{(k)}_a\theta^{(k)}_b,\qquad \theta^{(k)}_a\propto e^{h^{(k)}_a},
\end{equation}
a \emph{product of marginals}: the class carries per-class composition but no
residue--residue coupling, and \eqref{eq:generator} reduces to two independent single-site
chains (the spectator index $c$ drops out of the rate). This is exactly the additive
``$h\oplus h$'' \Potts\ matrix, and it is the null against which coupling is measured.
The coupling directions are the departures of $J$ from this additive subspace.

\section{The null-spiked, $L_1$-sparse prior}
\label{sec:prior}

We place a \emph{spike-and-slab} prior on each class's coupling, at the class level (the
``this class can be a pure null'' decision) with a Laplace ($L_1$) slab (within-class
sparsity). Introduce a class null-indicator $\gamma_k\in\{0,1\}$,
\begin{equation}
\label{eq:prior}
\gamma_k\sim\mathrm{Bernoulli}(\rho),\qquad
p\big(J^{(k)}\mid\gamma_k\big)=
\begin{cases}
\delta_0\!\big(J^{(k)}\big) & \gamma_k=0\ \ (\text{null}),\\[4pt]
\displaystyle\prod_{a<b}\tfrac{\lambda}{2}\,e^{-\lambda\,|J^{(k)}_{ab}|} & \gamma_k=1\ \ (\text{coupled}),
\end{cases}
\end{equation}
so that, marginalising $\gamma_k$, the prior on the coupling is the \emph{zero-adjusted}
(null-enriched) mixture
\begin{equation}
\label{eq:zeroadj}
p\big(J^{(k)}\big)=(1-\rho)\,\delta_0\!\big(J^{(k)}\big)
   +\rho\prod_{a<b}\tfrac{\lambda}{2}\,e^{-\lambda\,|J^{(k)}_{ab}|}
\end{equation}
--- an atom of mass $1-\rho$ on the product-of-marginals null, plus a sparse continuous
slab. The fields carry a smooth prior $h^{(k)}\sim\mathcal N(0,\sigma_h^2 I)$ (or a flat
one; composition is always present). Two hyperparameters control the two regularisers:
$\rho$ (the prior probability that a class couples at all --- ``null-enrichment'', small
$\rho$ demanding strong evidence before a class turns on) and $\lambda$ (the $L_1$ rate,
sparsifying \emph{which} residue pairs couple within an active class). Both admit
conjugate-ish hyperpriors, $\rho\sim\mathrm{Beta}(a_\rho,b_\rho)$ and a Gamma prior on
$\lambda$, so the corpus sets its own null-enrichment and sparsity level. An entry-level
variant replaces the single $\gamma_k$ by per-pair indicators $\gamma_{k,ab}$; the
class-level form is the one that answers ``is this an interaction class or a composition
class'', which is the quantity of interest.

\paragraph{Structural alternative.} Equivalently one may drop $\gamma_k$ and impose the
$L_1$ slab alone: \eqref{eq:zeroadj} with $\rho{=}1$ still yields exact zeros at the
optimum (the Laplace kink), so pure $L_1$ already gives within-class sparsity and, for
large $\lambda$, whole-class collapse to the null. The explicit spike is worth keeping
because it (i) separates the two questions --- \emph{whether} to couple ($\gamma_k$) from
\emph{how sparsely} ($\lambda$) --- and (ii) gives a calibrated posterior
$p(\gamma_k{=}1\mid\text{data})$, the phylogenetic-support probability we ultimately read
off.

\section{The MAP objective}
\label{sec:objective}

Fix the mixture E-step: responsibilities $r_{g,k,m}$ over class $k$ and rate bin $m$ for
each contact cluster $g$, and from them the responsibility-weighted, rate-aggregated
Holmes--Rubin sufficient statistics of class $k$ --- usage $N^{(k)}_{xy}$ (expected
$x\!\to\!y$ transitions) and dwell $T^{(k)}_x$ over the $A^2$ pair states. The expected
complete-data log-likelihood of class $k$ is the exponential family
\begin{equation}
\label{eq:ell}
\ell_k(h,J)=\sum_{x\ne y}N^{(k)}_{xy}\log Q_k(x\!\to\!y)
   -\sum_x T^{(k)}_x\!\!\sum_{y\ne x}\!Q_k(x\!\to\!y),
\end{equation}
a smooth concave-in-$\log Q$ function of $(h,J)$ through
\eqref{eq:generator}. The class MAP objective adds the prior:
\begin{equation}
\label{eq:map}
\mathcal F_k(h,J,\gamma_k)=\ell_k(h,J)-\tfrac{1}{2\sigma_h^2}\|h\|^2
   +\gamma_k\Big[\log\tfrac{\rho}{1-\rho}
   -\lambda\!\sum_{a<b}\!|J_{ab}|+\tbinom{A}{2}\log\tfrac{\lambda}{2}\Big]
   -(1-\gamma_k)\,\mathbb 1[J\ne0]\cdot\infty .
\end{equation}
Maximising over $(h,J)$ for each $\gamma_k$ and then over $\gamma_k$ gives the joint MAP;
alternating this per-class update with the standard mixture-weight and $\Gamma$+I updates
is a generalised EM.

\paragraph{The coupling gradient is a pair-flux balance.} Because $\log Z_k$ cancels,
$\partial\log Q_k[(a,c)\!\to\!(b,c)]/\partial J_{pq}=\tfrac12(\mathbb 1_{\{b,c\}=\{p,q\}}
-\mathbb 1_{\{a,c\}=\{p,q\}})$, so
\begin{equation}
\label{eq:grad}
\frac{\partial\ell_k}{\partial J_{pq}}
=\tfrac12\Big[\underbrace{U^{\to}_{pq}-U^{\leftarrow}_{pq}}_{\text{usage into}-\text{out of }(p,q)}
-\underbrace{\big(D^{\to}_{pq}-D^{\leftarrow}_{pq}\big)}_{\text{exposure-weighted rate flux}}\Big],
\end{equation}
where $U^{\to}_{pq}=\sum_{a\ne p}N^{(k)}_{(a,q),(p,q)}+\sum_{c\ne q}N^{(k)}_{(p,c),(p,q)}$
counts observed single-site moves that \emph{create} pair $(p,q)$ (and $U^{\leftarrow}$
those that destroy it), while $D^{\to}_{pq}=\sum_{a\ne p}T^{(k)}_{(a,q)}Q_k[(a,q)\!\to\!(p,q)]
+\dots$ is the model's expected creation rate weighted by dwell. At the unpenalised
optimum $\partial\ell_k/\partial J=0$ is a detailed pair-flux balance: the observed net
flux into each pair matches the model's. This is the phylogenetic analogue of the
plmDCA moment-matching condition (observed pair frequency $=$ model pair frequency), with
the tree, branch lengths, and rate heterogeneity carried exactly through
$N^{(k)},T^{(k)}$.

\section{Inference}
\label{sec:inference}

\paragraph{Coupled M-step ($\gamma_k{=}1$): $L_1$-proximal ascent.} Maximise
$\ell_k(h,J)-\lambda\sum_{a<b}|J_{ab}|$ by proximal gradient: a smooth step on
$(h,J)$ using \eqref{eq:grad}, followed by soft-thresholding of the coupling,
\begin{equation}
\label{eq:soft}
J_{ab}\leftarrow\soft_{\eta\lambda}\!\Big(J_{ab}+\eta\,\tfrac{\partial\ell_k}{\partial J_{ab}}\Big),
\qquad \soft_{\tau}(x)=\mathrm{sign}(x)\max(|x|-\tau,0),
\end{equation}
with step $\eta$ (a diagonal/backtracking choice keeps it stable given the pair-flux
Hessian). The threshold zeroes every pair whose net flux imbalance is below $\lambda$, so
an active class keeps only the residue pairs it genuinely needs --- e.g.\ the acid/base
salt-bridge cells for the charge class. The fields update in closed form as usual
($h\!\propto$ log-composition plus the Gaussian shrinkage).

\paragraph{Null branch ($\gamma_k{=}0$): $J{=}0$.} Fit $h$ alone; the class is a
product-of-marginals composition class, scored by $\ell_k(h,0)$.

\paragraph{The null decision.} Choose $\gamma_k$ by the (penalised) evidence contrast
\begin{equation}
\label{eq:decision}
\gamma_k=1\iff
\underbrace{\big[\ell_k(\hat h_1,\hat J)-\lambda\!\sum_{a<b}\!|\hat J_{ab}|\big]}_{\text{coupled fit}}
-\underbrace{\ell_k(\hat h_0,0)}_{\text{null fit}}
\;>\;\log\tfrac{1-\rho}{\rho}-\tbinom{A}{2}\log\tfrac{\lambda}{2},
\end{equation}
i.e.\ the coupling turns on only when its likelihood gain clears the log-odds of the
null prior plus the slab's normalisation. A soft (variational/EM) treatment replaces the
hard rule by the responsibility
$q(\gamma_k{=}1)=\sigma\big(\text{LHS}-\text{RHS of \eqref{eq:decision}}\big)$
and marginalises it in the class-mixing update; a Laplace approximation of the slab
evidence around $\hat J$ sharpens the RHS. Either way \eqref{eq:decision} is a
Bayes-factor test for ``does this class carry phylogenetically-supported coupling''.

\paragraph{Hyperparameters.} With $\rho\sim\mathrm{Beta}(a_\rho,b_\rho)$, the conditional
update is $\hat\rho=(a_\rho-1+\sum_k q(\gamma_k))/(a_\rho+b_\rho-2+K)$; $\lambda$ updates
from its Gamma prior against $\sum_k\gamma_k\sum_{a<b}|J^{(k)}_{ab}|$. Both are cheap and
make the null-enrichment and sparsity data-driven.

\section{The read-off: coupling that survives as phylogenetic signal}
\label{sec:readoff}

The point of the prior is what it lets us report. After fitting:
\begin{itemize}
\item $\sum_k q(\gamma_k{=}1)$ --- the expected number of \emph{interaction} classes (as
opposed to composition classes); with strong null-enrichment we expect this to be small,
consistent with the $\sim$$0.008$-nat coupling gain.
\item Per class, the surviving support $\{(a,b):\hat J^{(k)}_{ab}\ne0\}$ --- the residue
pairs the data keep after $L_1$; the charge/salt-bridge cells should dominate, matching
the sign-consistent charge axis that survives the washout (paper2b, \S\,axes).
\item The total selected coupling mass $\sum_k w_k\,q(\gamma_k)\,\|\hat J^{(k)}\|_1$ and the
held-out likelihood gain over the all-null ($\gamma\equiv0$) fit --- a single, honest
number for ``how much coupling is real'', now defended by a model that was free to keep
none.
\end{itemize}
Because the null is nested ($\gamma\equiv0$ is the product-of-marginals mixture) and the
prior is free to select it, the spiked model cannot do worse than the null on the training
objective and reports coupling only where it is paid for --- exactly the discipline the
small measured coupling gain calls for.

\section{Relation to prior work and to the rest of the paper}

The slab is the phylogenetic counterpart of sparse-$L_1$ DCA (plmDCA), with the static
pseudolikelihood replaced by the exact tree/branch/rate likelihood \eqref{eq:ell} and the
partition-function cancellation of \eqref{eq:generator}; the spike adds the
group-sparse/zero-inflated ``is this class an interaction class'' decision that a flat DCA
prior lacks. Against the paper's models: the all-null limit ($\gamma\equiv0$) is the
mixture of products-of-marginals (the additive-\Potts\ null of \S\,axes); switching every
$\gamma_k{=}1$ with $\lambda{=}0$ recovers the unpenalised Metropolis--\Potts\ mixture of
\texttt{tab:pairfit}; and the spiked prior interpolates, selecting per class. It replaces
the abandoned lumpable mixture entirely: rather than constraining the dynamics (which, as
shown, forces coupling out of the stationary), it constrains the \emph{prior over
couplings}, which is where the small surviving signal should be adjudicated.

\end{document}
